% Pursuing Stacks (cotes 134-2 to 134-8): archive cote 134-5, site page 108 (PDF page 109), typescript page 295.
% typescript. Draft compared: ps5.tex:189-246.
% AI-transcribed by Claude Opus 5.5 (claude-opus-5-5, adaptive thinking, effort high), Message Batches API, batch ps2 (msgbatch_01NVwr1Yuo8qJf5rykeGBXPK), 2026-11-12; not yet checked by eye.
\scanpage{134-5/108}
\mspage{295}

cylinder of $0\to X$ ),  thus the \xout{\illegible{1}} set of aspheric objects follows
formally from the functor $H^{o}$  (as the objects such that $H^{i}(X) = 0$ for all
$i$), much as in the non commutative set-up the functor $\hand{\pi}_{o}$ for a contractibi-%
lity structure determines the latter,  and hence also the corresponding
asphericity structure.

Thus,  the question arises of formulating the basic structures,
namely \xout{\unsure{pure}} contractibility structures and asphericity structures,  in an
autodual way,applying both to the autodual situation just described,  and
to the non autodual one we have been working out in the notes  -  and if
possible even,  \over{\unsure{in}}{to} all \xout{\illegible{1}} situations met with so far where a homology
or homotopy \inserted{theory} of some kind or other has turned up.  Of course,  it is by no
means sure a priori that we can do so,  by keeping first nicely apart the
two sets of notions (contractibility and asphericity) namely defining
them separately,  and then showing that a contractibility structure de-%
termines an asphericity structure\hand{,} and is determined by the latter.  Maybe
we'll have to define from the outset a richer kind of structure,  were
both ``pure'' homotopy notions and asphericity notions are involved.  Also,
the familiar generalization of mapping cylinders,  namely homotopy
fibers and cofibers,  and the corresponding long exact sequences, will
evidently play an important role in the structure to be described.  Now
this again ties in with the corresponding structure of a derivator,  as
contemplated in section 69,  namely ``integration'' and ``cointegration'' of
diagrams in the given category.

Definitely,  this ref\over{f}{l}ection is \xout{\unsure{leading to}}\inserted[typed]{going to} lead well beyond the
scope of the present part IV -  it relates rather to part C) of the
working program envisioned by the end of Mai (section 71,  p.207-210),
\margin{$\times$} and rather belongs to \uline{part VI} of these notes,  which presumably will cen-%
ter around the notion of a derivator.  There is  however a more technical
question, \inserted{and} of more limited scope,  which deserves some thought and goes
somewhat in the same direction,  namely:  how to define (for a given
contractibility or asphericity structure on M) an \uline{induced} structure on
a category $M_{/a}$ ,  where a is in M ?   There is a little perplexity in
my mind,  even when M is of the type $C^{\wedge}$ say\hand{,} with C a contractor\hand{,} and a \xout{\unsure{is}}
in C ,  taking the canonical structures on $C^{\wedge}$ -  because with the most
evident choice of an``induced''  asphericity structure on $(C^{\wedge})_{/a} \hand{\simeq} (C_{/a})^{\wedge}$,
\xout{this structure} namely the usual notion of aspheric objects,  this structu\inserted{re}

\margin{108}
